\chapter{RESULTS AND DISCUSSION}

\section{SANITY CHECK: OVERFITTING ON 10 SAMPLES}

Before full training, a sanity check verified that each neural network model
could memorise a tiny batch of 10 training examples. With a plain (unweighted)
BCE loss and a learning rate of $5 \times 10^{-3}$, both the CNN and TCN reduced
their training loss below $10^{-2}$ within 200 epochs. This confirms that the
forward and backward passes are implemented correctly and that the architecture
has sufficient capacity for the task.

\section{TRAINING CONVERGENCE}

\subsection{CNN}

The CNN was trained for 5 epochs before early stopping halted training. The best
checkpoint was saved at \textbf{epoch 3} with validation F1\,=\,0.9974 and
validation MATE\,=\,0.246 samples. The rapid convergence is consistent with the
relatively large model capacity (30\,M parameters) and the high information
density of the weighted loss.

% TODO: Insert training curve figure here.
% \begin{figure}[ht]
%     \centering
%     \includegraphics[width=\linewidth]{images/cnn_training_curves.pdf}
%     \caption{CNN training and validation loss (left) and validation F1-score
%              (right) over epochs. The red dashed line marks the best epoch.}
%     \label{fig:cnn_curves}
% \end{figure}

\subsection{TCN}

The TCN was trained for 33 epochs. The best checkpoint was saved at
\textbf{epoch 23} with validation F1\,=\,0.9982 and validation
MATE\,=\,0.237 samples. The learning rate was reduced automatically twice during
training (ReduceLROnPlateau), which produced the characteristic step-wise
improvement visible in the validation curve.

% TODO: Insert training curve figure here.
% \begin{figure}[ht]
%     \centering
%     \includegraphics[width=\linewidth]{images/tcn_training_curves.pdf}
%     \caption{TCN training and validation loss (left) and validation F1-score
%              (right) over epochs. The red dashed line marks the best epoch.}
%     \label{fig:tcn_curves}
% \end{figure}

\section{QUANTITATIVE COMPARISON}

Table~\ref{tab:results} summarises the performance of all three methods on the
validation set using the best saved checkpoint for each neural network model.

\begin{table}[ht]
    \centering
    \caption{Axle-level detection performance on the validation set.
             MATE is the Mean Absolute Timing Error in samples.
             $\pm$5 sample tolerance window, greedy matching.}
    \label{tab:results}
    \begin{tabular}{lccccr}
        \toprule
        \textbf{Method} & \textbf{Precision} & \textbf{Recall}
                        & \textbf{F1} & \textbf{MATE} & \textbf{Params} \\
        \midrule
        Baseline (SG + peaks) & -- & -- & -- & -- & 0 \\
        % TODO: Fill in baseline numbers after running notebook 02.
        CNN (U-Net, epoch 3)  & 0.9954 & 0.9995 & 0.9974 & 0.246 & $\sim$30\,M \\
        TCN (epoch 23)        & \textbf{0.9965} & \textbf{0.9999}
                              & \textbf{0.9982} & \textbf{0.237} & $\sim$2.7\,M \\
        \bottomrule
    \end{tabular}
\end{table}

Both deep learning models achieve near-perfect axle detection. The TCN
outperforms the CNN on all four metrics while using 11 times fewer parameters.
The MATE of 0.24 samples corresponds to approximately 2.4\,ms at a typical
B-WIM sampling rate of 100\,Hz, which is operationally negligible for weight
calculation purposes.

\section{VISUAL INSPECTION}

Figure~\ref{fig:visual} shows the TCN's predictions on six randomly selected
test-set records. In each panel, the raw strain signal (blue) is overlaid with
the scaled TCN output probability (orange). Ground-truth axle positions are
marked by green vertical lines and predicted positions by red dashed lines.

Across all six examples the TCN correctly identifies the axle count and places
predictions within 1--2 samples of the ground truth. No false positives are
observed in these examples, and the probability output forms a sharp, narrow peak
above each axle event, confirming that the model has learned a precise temporal
localisation rather than a diffuse detection.

% TODO: Insert visual inspection figure here (generated from notebook 04, cell 14).
% \begin{figure}[ht]
%     \centering
%     \includegraphics[width=\linewidth]{images/tcn_visual_inspection.pdf}
%     \caption{TCN predictions on six test-set records. Blue = strain signal;
%              orange = TCN output probability (scaled to signal amplitude range);
%              green lines = ground-truth axle positions;
%              red dashed lines = predicted axle positions.}
%     \label{fig:visual}
% \end{figure}

\section{DISCUSSION}

\subsection{Architecture vs.\ Model Size}

The most striking finding is that the TCN, with 2.7\,M parameters, consistently
outperforms the CNN, which has 30\,M parameters. This supports the hypothesis
that the key design choice for this task is receptive field coverage rather than
raw model capacity. The TCN's dilated blocks ensure that every output position
can attend to the full 1,300-sample signal, while the CNN's skip connections
compensate for the loss of temporal resolution in the encoder but cannot replicate
this global context.

\subsection{Convergence Speed}

The CNN converged in only 3 effective epochs, whereas the TCN required 23. This
likely reflects the CNN's larger parameter count providing more initial
flexibility, but also a tendency to overfit sooner. The slower, more regular
convergence of the TCN --- with two visible learning rate reduction steps ---
suggests better regularisation and a smoother optimisation landscape.

\subsection{Practical Implications}

A MATE of 0.24 samples translates to sub-millisecond timing accuracy at 100\,Hz,
well within the tolerance required for accurate weight calculation. Both models
achieve near-perfect recall ($\geq 0.9995$), meaning very few axles are missed.
Precision is slightly lower ($\sim 0.9954$--$0.9965$), indicating occasional
false positives, but these remain rare (fewer than 1 in 200 detected axles).

\subsection{Limitations}

\begin{itemize}
    \item \textbf{Record-level data split}: The train/validation/test split is
          performed by vehicle record, not by day or week. Crossings from the same
          day may share environmental conditions (temperature, noise floor),
          potentially inflating generalisation estimates. A chronological
          (day-level) split would provide a stricter evaluation.
    \item \textbf{Single installation}: All data originates from one bridge and
          sensor configuration. Performance on a different bridge or sensor type
          is unknown.
    \item \textbf{Raised axles}: Vehicles with a raised (lifted) axle produce a
          zero or near-zero signal for that axle. Neither model was explicitly
          tested on such cases.
    \item \textbf{Robustness to noise and speed variation}: A systematic
          robustness experiment (e.g., adding synthetic Gaussian noise or
          time-stretching signals to simulate speed variation) was outside the
          scope of this work.
\end{itemize}